\chapter{Zusammenfassung Fachdidaktik Geographie}
\section{Leistungsbeurteilung und Prüfungsziele}
\begin{itemize}[leftmargin=*]
  \item Kenntnis kantonaler und ausserkantonaler Lehrpläne
  \item Didaktische und methodische Zugänge zu sozial- und naturwissenschaftlichen Themen
  \item Einbezug von Präkonzepten und Alltagserfahrungen der Lernenden
  \item Analyse und didaktische Rekonstruktion komplexer Inhalte
  \item Grundfertigkeiten: Beobachten, beschreiben, analysieren, vernetzen, experimentieren etc.
  \item Unterrichtskritik und Weiterentwicklung
  \item Ziel- und Kompetenzorientierung im eigenen Unterricht
  \item Durchführung von Unterricht an verschiedenen Lernorten
\end{itemize}

\section{Unterrichtsprinzipien im Geografieunterricht}
\subsection{Überblick}
\begin{itemize}[leftmargin=*]
  \item Problemorientierung
  \item Zielorientierung
  \item Schülerorientierung
  \item Inhaltliche Strukturierung
  \item Wissenschaftsorientierung
  \item Exemplarisches Prinzip
  \item Anschauung und Raumbezug
  \item Aktualitätsprinzip
\end{itemize}

\subsection{Problemorientierung}
\begin{itemize}[leftmargin=*]
  \item Relevante Problemstellungen fördern aktives, vernetztes Denken
  \item Beispiel: Tourismus in Interlaken – Wer profitiert?
\end{itemize}

\subsection{Zielorientierung}
\begin{itemize}[leftmargin=*]
  \item Steuerung aller Unterrichtsebenen
  \item Lehrplan- und kompetenzbasierte Planung
\end{itemize}

\subsection{Schülerorientierung}
\begin{itemize}[leftmargin=*]
  \item Einbezug von Vorwissen, Alltagserfahrung, Interessen
  \item Differenzierung und Förderung der Selbstwirksamkeit
\end{itemize}

\subsection{Inhaltliche Strukturierung}
\begin{itemize}[leftmargin=*]
  \item Transparente Gliederung des Fachinhalts
  \item Ermöglicht Transfer und Einbettung in bestehende Wissensstrukturen
\end{itemize}

\subsection{Wissenschaftsorientierung}
\begin{itemize}[leftmargin=*]
  \item Aktueller Forschungsstand, sachlogische Strukturierung, Modelle
  \item Methodenkompetenz und fachliche Korrektheit
\end{itemize}

\subsection{Exemplarisches Prinzip}
\begin{itemize}[leftmargin=*]
  \item Auswahl von Lerninhalten mit übertragbaren Einsichten
  \item Fokus auf Elementares und Fundamentales
\end{itemize}

\section{Kompetenzorientierter Unterricht}
\begin{itemize}[leftmargin=*]
  \item Fokussierende Lernaufgaben (AEL)
  \item Conceptual Change und kognitive Aktivierung
  \item Transfer fördern, z.B. mithilfe von Erklärungsansätzen
\end{itemize}

\section{Bildung für nachhaltige Entwicklung (BNE)}
\begin{itemize}[leftmargin=*]
  \item Leitprinzip der Schulgeografie
  \item Integration der Dimensionen Umwelt, Gesellschaft, Wirtschaft
  \item Beitrag zu den SDGs (z.B. SDG 4)
  \item Starke/schwache NE, didaktische Umsetzung anhand z.B. Strategiespiel \emph{Bougouni}
\end{itemize}

\section{Methoden und Medien}
\begin{itemize}[leftmargin=*]
  \item Frontalunterricht, Mystery, Gruppenpuzzle, Diskussion
  \item Medienverbund mit Bildern, Filmen, Karten, Simulationen
  \item Kritische Reflexion medialer Inhalte
\end{itemize}

\section{Planung und Unterrichtsphasen}
\begin{itemize}[leftmargin=*]
  \item Modell der didaktischen Rekonstruktion
  \item Phasen: Einstieg, Erarbeitung, Sicherung, Anwendung, Transfer
\end{itemize}